\documentclass[11pt,a4paper]{article}
\usepackage[T1]{fontenc}
\usepackage[utf8]{inputenc}
\usepackage{lmodern}
\usepackage[margin=2.5cm]{geometry}
\usepackage{booktabs}
\usepackage{tabularx}
\usepackage{longtable}
\usepackage{enumitem}
\usepackage{listings}
\usepackage{xcolor}
\usepackage[hidelinks]{hyperref}

\lstset{
  basicstyle=\ttfamily\small,
  columns=fullflexible,
  keepspaces=true,
  frame=single,
  rulecolor=\color{gray!50},
  backgroundcolor=\color{gray!8},
  breaklines=true,
  upquote=true,
  showstringspaces=false,
  extendedchars=true,
  inputencoding=utf8,
  literate={é}{{\'e}}1 {è}{{\`e}}1 {à}{{\`a}}1 {ç}{{\c c}}1 {É}{{\'E}}1,
  xleftmargin=0.5em,
  xrightmargin=0.5em
}
\setlength{\emergencystretch}{3em}
\newcommand{\kbd}[1]{\texttt{#1}}
\newcolumntype{L}{>{\raggedright\arraybackslash}X}

\title{LeRobot (Python) User Manual}
\author{}
\date{9 November 2026}

\begin{document}
\maketitle

LeRobot is a drawing robot you steer on screen: either live, with the mouse
and keyboard (Dessiner), or by writing a short Python program of robot
commands (such as \kbd{mon\_dessin.py}). It is the Python translation of
the OCaml version, itself a translation of the 1998 Caml Light version. It
lives in the folder \kbd{LeRobot\_Python\_FR} of the \kbd{LeRobot} folder,
which can be anywhere (commands below are typed in a Terminal window
opened in the \kbd{LeRobot} folder); the
installation is described in the \emph{LeRobot (Python) Installation
Manual} (\kbd{LeRobot\_Install.pdf}).

\tableofcontents

%======================================================================
\section{Getting started}

You need a Python linked to Tk 8.6 or later, for example the one from
python.org (see the installation manual). There is nothing to compile and
no XQuartz: the robot window is an ordinary macOS window.

There are three ways to start a program:
\begin{itemize}
\item \textbf{The application} \kbd{LeRobotPython.app} (in
  \kbd{LeRobot\_Python\_FR}): double-click it to run Dessiner; drop a \kbd{.py}
  file on its icon to run that program.
\item \textbf{The script} \kbd{robot.sh}, in a Terminal window opened in
  \kbd{LeRobot\_Python\_FR}: \kbd{./robot.sh path/to/Program.py}. The program
  can be in any folder.
\item \textbf{Python directly}, for the files of the \kbd{LeRobot\_Python\_FR}
  folder itself: \kbd{python3 Dessiner.py}, \kbd{python3 mon\_dessin.py}.
\end{itemize}
If the \kbd{python3} of the Terminal is the one of macOS (linked to the old
Tk 8.5), LeRobot notices it and starts itself again with the python.org
Python; nothing has to be done.

The window is $800 \times 700$ points. The large white area is where the
robot draws; the strip along the bottom is the message bar.

%======================================================================
\section{Part 1 --- Interactive drawing (Dessiner)}

In Dessiner you give the robot commands with single keys and move it with
the mouse or the arrow keys; press Enter when you are done, then \kbd{y} to
save the drawing as a program.

\subsection{Launch}
Double-click \kbd{LeRobotPython.app}, or:
\begin{lstlisting}[language=bash]
cd LeRobot_Python_FR && python3 Dessiner.py
\end{lstlisting}
The robot appears in the centre, pointing north. The window comes to the
front with the keyboard focus; if keys do nothing, click in the window.

\subsection{The window}
\begin{itemize}
\item \textbf{Drawing area} --- the large white area.
\item \textbf{Message bar} (bottom) --- a reminder of the keys, ending with
  \kbd{* :cc w :peindre\ \ Entree :fin}.
\item \textbf{Colour box} (small robot near the right of the bar) --- shows
  the current pen colour.
\item \textbf{Enter icon} (far right) --- the Enter symbol; it blinks when
  the robot waits for a key.
\end{itemize}

\subsection{Moving the robot with the mouse or the arrow keys}
Click on a grid point (or move the pointer with the button held down) and
the robot jumps to the nearest grid point, without drawing. The arrow keys
move it one grid step left, right, up or down, also without drawing. When
the pointer moves with the button released, the robot does not move.

If the robot would go past the edge of the drawing area, it stays where it
is and the terminal prints \kbd{** SORTIE INTERDITE ** translater}
(``forbidden exit''), with a beep-beep. This is only a warning.

\subsection{Keyboard commands}
The robot draws only when its pen is down (\kbd{b}; it is down at the
start). One step is one grid square.

\begin{center}
\begin{tabularx}{\linewidth}{@{}l l L@{}}
\toprule
\textbf{Key} & \textbf{Command} & \textbf{What it does} \\
\midrule
\kbd{b} & bc & Pen down: the robot draws as it moves \\
\kbd{l} & lc & Pen up: the robot moves without drawing \\
\kbd{a} & av & Move forward one step \\
\kbd{)} & vd & Curved turn to the right (draws an arc) \\
\kbd{(} & vg & Curved turn to the left (draws an arc) \\
\kbd{+} & pqtd & Quarter turn right ($90^\circ$) \\
\kbd{-} & pqtg & Quarter turn left ($90^\circ$) \\
\kbd{d} & phtd & Eighth turn right ($45^\circ$) \\
\kbd{g} & phtg & Eighth turn left ($45^\circ$) \\
\kbd{c} & dc & Draw a small cross (+) \\
\kbd{p} & dp & Draw a point \\
\kbd{.} & dg & Draw the grid points \\
\kbd{*} & cc & Next pen colour (shown in the colour box) \\
\kbd{w} & peindre & Paint the closed area behind the robot with the pen colour \\
Arrow keys & translater & Move one step left, right, up or down, without drawing \\
Backspace & --- & Undo the last command \\
Enter & --- & Finish the drawing \\
\bottomrule
\end{tabularx}
\end{center}

Colours cycle in this order: white, yellow, blue, violet, green, red,
indigo, black. To paint, put the robot so that its back is inside a closed
area, then press \kbd{w}.

\subsection{Finishing and saving}
\begin{enumerate}
\item Press \textbf{Enter}. The bar asks:
  \emph{Voulez-vous sauver dans mon\_dessin.py ? (y/n)}
\item Press \kbd{y} to save. Every command you used, including mouse and
  arrow moves (written as \kbd{translater(dx, dy)}) and colour changes
  (written as \kbd{cc(n)} with the colour name in a comment), is written to
  \kbd{mon\_dessin.py} as a Python program (see Part~2), and the window
  closes.
\item Press any other key to close without saving.
\end{enumerate}
Saving always overwrites \kbd{mon\_dessin.py}, in the folder of
Dessiner. To keep an earlier drawing, copy that file under another name
first.

%======================================================================
\section{Part 2 --- Running a robot program}

A robot program is a Python file that lists robot commands in order. It
draws the picture by itself, then waits so you can replay it at another
speed or close it with Esc.

\subsection{Run mon\_dessin.py}
\begin{lstlisting}[language=bash]
cd LeRobot_Python_FR && python3 mon_dessin.py
\end{lstlisting}
or \kbd{make mon\_dessin}, or drop \kbd{mon\_dessin.py} on
\kbd{LeRobotPython.app}. A saved drawing looks like this:
\begin{lstlisting}[language=Python]
# Mon joli dessin
from LeRobot import *

init_pict(); ce()
av(); av(); pqtd(); av(); av(); cc(2)  # bleu
av(); translater(5, 3); av(); vg(); av(); st()
\end{lstlisting}

\subsection{Write your own program}
Create a file anywhere, for example \kbd{maison.py}:
\begin{lstlisting}[language=Python]
from LibRobot import *

# Dessiner un carre
av(); av(); pqtd(); av(); av(); pqtd()
av(); av(); pqtd(); av(); av(); pqtd()
st()
\end{lstlisting}
Run it with \kbd{./robot.sh path/to/maison.py} from the
\kbd{LeRobot\_Python\_FR} folder, or drop it on \kbd{LeRobotPython.app}.

Every program follows the same pattern:
\begin{itemize}
\item \kbd{from LibRobot import *} first. It opens the window, puts the
  robot in the centre, pointing north, pen down, in black, with a grid of
  12 points (the settings of the book's programs). \kbd{from LeRobot import *}
  gives the original interface used by \kbd{mon\_dessin.py}: a grid of 30
  points, and the program must start with \kbd{init\_pict(); ce()}.
\item The robot starts with its pen down, so every move draws: call
  \kbd{lc()} to move without drawing, and \kbd{bc()} to draw again.
\item \kbd{st()} last: shows the speed screen and keeps the window open.
  Without it, the bar says \emph{Fin du programme : taper une touche ou
  cliquer} and the window closes at the next key or click.
\item Commands without arguments take \kbd{()}; commands are separated by
  \kbd{;} or written on separate lines.
\end{itemize}

\subsection{Programs of the book}
The folder \kbd{livre} holds the programs of the book \emph{Premières
leçons de programmation}, one subfolder per chapter:
\begin{lstlisting}[language=bash]
./robot.sh livre/20_grille_polaire/LimaconDePascal.py
make run PROG=livre/21_recursivite/CourbeDuDragon
make verifier                      # all programs, one after the other
make verifier DIR=02_interaction   # one chapter
make verifier DEPUIS=FriseDeTriangles   # resume at a program
\end{lstlisting}
With \kbd{make verifier}, answer in the robot window when a drawing is
finished: Enter = correct (next program), \kbd{r} = run it again, \kbd{q}
or Esc = stop (the command to resume is printed), any other key starts a
remark, ended by Enter, noted in \kbd{verification.txt}. A program that
never ends by itself (\kbd{Boucle1}, \kbd{Boucle2}) is stopped with Esc.
Programs that ask questions show them in the bottom bar and in the
Terminal; answer in either. \kbd{LEROBOT\_MUET=1} turns the sound off.

\subsection{Replaying at another speed}\label{sec:replay}
When the program reaches \kbd{st()}, the bar shows:
\begin{lstlisting}
+ : plus vite   - : moins vite   vitesse : N     autre touche : redessiner   Echap : quitter
\end{lstlisting}
\begin{itemize}
\item \kbd{+} / \kbd{-} --- change the speed from 0 to 10 (default 10, the
  fastest). Only the number changes.
\item Any other key, for example the space bar --- clear the window and
  redraw at the chosen speed; the speed screen comes back when it is done.
\item Esc --- close the window.
\end{itemize}
At speed 1 each command takes about 0.1\,s. Speed 0 is step-by-step mode:
the robot does one command each time you press a key. While the robot
draws, Esc (or \kbd{\S} or \kbd{\#}) stops the program.

\subsection{Command reference}
One step is one grid square. Commands are in French; the table gives the
English meaning. Commands with two numbers take them as two arguments:
\kbd{deplacer(x, y)}.

\begin{longtable}{@{}p{0.42\linewidth} p{0.53\linewidth}@{}}
\toprule
\textbf{Command} & \textbf{Meaning} \\
\midrule
\endhead
\bottomrule
\endfoot
\kbd{av()} & Move forward one step \\
\kbd{avf(n)} & Move forward $n$ steps (backwards if $n < 0$) \\
\kbd{vd()} / \kbd{vg()} & Curved turn right / left (draws an arc) \\
\kbd{phtd()} / \kbd{phtg()} & Turn $45^\circ$ right / left \\
\kbd{pqtd()} / \kbd{pqtg()} & Turn $90^\circ$ right / left \\
\kbd{p3htd()} / \kbd{p3htg()} & Turn $135^\circ$ right / left \\
\kbd{pdt()} & Turn around ($180^\circ$) \\
\kbd{pn()} \kbd{pne()} \kbd{pe()} \kbd{pse()} \kbd{ps()} \kbd{psw()} \kbd{pw()} \kbd{pnw()} & Face N, NE, E, SE, S, SW, W, NW \\
\kbd{pivoter(n)} & Face direction $n$ (0 = N, 1 = NE, 2 = E \dots\ 7 = NW) \\
\kbd{rt(a)} & Turn by $a$ degrees (anticlockwise), keeping the step length \\
\kbd{bc()} / \kbd{lc()} & Pen down / pen up \\
\kbd{dp()} / \kbd{dc()} & Draw a point (a square of side \kbd{ec}) / a small cross \\
\kbd{cc(n)} & Pen colour $n$ (0 white, 1 yellow, 2 blue, 3 violet, 4 green, 5 red, 6 indigo, 7 black; also \kbd{blanc} \dots\ \kbd{noir} with LibRobot) \\
\kbd{ec(n)} & Pen thickness $n$ (a square pen, as on the Macintosh) \\
\kbd{cf(n)} & Clear the window and fill it with background colour $n$ \\
\kbd{peindre()} & Paint the closed area behind the robot with the pen colour \\
\kbd{translater(dx, dy)} & Move $dx$ steps across and $dy$ steps up, without drawing \\
\kbd{tracer(x, y)} & Go to window pixel $(x, y)$, drawing if the pen is down \\
\kbd{deplacer(x, y)} & Go to grid point $(x, y)$, counted from the centre \\
\kbd{tl(l)} & Move forward $l$ points, the grid following the robot \\
\kbd{ce()} & Go to the centre \\
\kbd{eh()} \kbd{eb()} \kbd{ad()} \kbd{ag()} & Go to the top, bottom, right or left edge \\
\kbd{dg()} & Draw the grid points \\
\kbd{lg(x)} / \kbd{lgX(x)} / \kbd{lgY(x)} & Set the step size (both directions / across / up) \\
\kbd{lgRT(r, t)} & Step of length $r$ at angle $t$ (polar grid) \\
\kbd{montrerRobot()} / \kbd{cacherRobot()} & Show / hide the robot \\
\kbd{autoriserSorties()} / \kbd{interdireSorties()} & Allow / forbid moving past the window edge \\
\kbd{vt(n)} & Set the speed (0 to 10) \\
\kbd{pp()} / \kbd{ex()} & Step-by-step mode / full speed \\
\kbd{message("text")} & Write a message in the bottom bar \\
\kbd{ecrire("text")} & Write a text at the robot's position (LibRobot) \\
\kbd{lireEntier()} / \kbd{lireCar()} & Read a number / a key (LibRobot) \\
\kbd{write(\dots)} / \kbd{writeln(\dots)} & Write in the Terminal (LibRobot) \\
\kbd{st()} & End: show the speed screen \\
\end{longtable}

Test functions such as \kbd{tb()} (would the next step leave the window?)
and the values \kbd{valX()}, \kbd{valLgX()}, \dots\ are defined in
\kbd{LeRobot.py}; the additions of the book (sounds,
reading, \kbd{interactionTP}, \dots) are in \kbd{LibRobot.py}. The
differences with the OCaml version are described in the \kbd{README.md}
file of \kbd{LeRobot\_Python\_FR}.

%======================================================================
\section{Troubleshooting}

\begin{longtable}{@{}p{0.27\linewidth} p{0.33\linewidth} p{0.34\linewidth}@{}}
\toprule
\textbf{Symptom} & \textbf{Cause} & \textbf{Fix} \\
\midrule
\endhead
\bottomrule
\endfoot
The program hangs, no window & No Python with Tk 8.6 on this Mac & Install Python from python.org (installation manual) \\
\addlinespace
\emph{No module named 'LibRobot'} & The program is not in \kbd{LeRobot\_Python\_FR} and was started with \kbd{python3} & Use \kbd{./robot.sh Program.py} or \kbd{LeRobotPython.app} \\
\addlinespace
\emph{Erreur ligne N de Prog.py : le robot sortirait du cadre (av)} & The robot was asked to move past the window edge & Change the program, or call \kbd{autoriserSorties()} \\
\addlinespace
\kbd{** SORTIE INTERDITE ** translater} in the terminal & In Dessiner, a move past the window edge & Only a warning \\
\addlinespace
The program never ends & An endless loop (\kbd{while True}, endless recursion) & Press Esc in the robot window \\
\addlinespace
Keys do nothing in the window & The window does not have keyboard focus & Click in the window \\
\addlinespace
Digits are not accepted & French AZERTY keyboard: the top row gives \& é " \dots\ without Shift & Questions (\kbd{lireEntier}) accept the top row; for \kbd{lireCar}, use Shift or the numeric keypad \\
\addlinespace
A saved drawing was lost & Saving from Dessiner always overwrites \kbd{mon\_dessin.py} & Copy \kbd{mon\_dessin.py} under a new name before saving again \\
\addlinespace
Text font cannot be changed & \kbd{policeTexte} and \kbd{styleTexte} have no effect & None; texts use Helvetica 13 \\
\end{longtable}

\end{document}
